%% notes-tex/028-gene-essentiality/tables/fold_ensembles.tex
%% GENERATED by experiments/028-gene-essentiality/scripts/notes_tex_tables.py
%% SOURCE: results/<prefix>_ensemble.json
\begin{table}[H]
  \centering
  \footnotesize
  \caption[Fold ensembles]{Every fold ensemble in the document on the 195 held-out genes: five models trained on the five cross-validation files (each a different fifth of the 787 labeled non-test genes held for validation), their per-gene logits averaged and scored once, Flux Cone Learning's own test protocol. Fold AUROCs are the single models on the held-out genes; the ensemble carries a 2,000-draw gene bootstrap. Written by \texttt{cv\_ensemble\_readout.py}. Bold marks the highest value in a column.}
  \label{tab:ensembles}
  \begin{tabular}{@{}p{62mm}rlrl@{}}
    \hline
    cell & folds & fold AUROCs & fold mean & ensemble [95\% CI] \\
    \hline
    round 3: pooled transformer, ProtT5, no readout & 5 & 0.780, 0.806, 0.797, 0.751, 0.789 & 0.785 & 0.795 [0.71, 0.88] \\
    round 5: ProtT5 + FBA affine & 5 & 0.814, 0.824, 0.839, 0.775, 0.853 & 0.821 & 0.822 [0.74, 0.90] \\
    round 12: chromosome-pathway indicator + ProtT5 + FBA affine (the reported number) & 5 & 0.860, 0.899, 0.889, 0.875, 0.884 & 0.881 & 0.893 [0.84, 0.94] \\
    round 17: chromosome-pathway indicator + ProtT5 + the widened affine (ratios, statistics, pFBA vector) & 5 & 0.863, 0.860, 0.887, 0.840, 0.882 & 0.867 & \textbf{0.901} [0.85, 0.95] \\
    \hline
  \end{tabular}
\end{table}
